\documentclass[11pt,a4paper]{article}
\usepackage[margin=2.5cm]{geometry}
\usepackage{amsmath,amssymb,amsthm}
\usepackage{booktabs}
\usepackage[hidelinks]{hyperref}

\newcommand{\C}{\mathbb{C}}
\newcommand{\B}{\mathcal{B}}
\newcommand{\G}{\mathcal{G}}
\newcommand{\J}{\mathcal{J}}
\newtheorem{lemma}{Lemma}

\title{Block encodings for numerical linear algebraists\\[2pt]
\large How the combination rules are realized\\
(supplement to ``Preconditioning for quantum linear solvers in PDE-constrained optimization'')}
\author{Martin Stoll \and John W. Pearson}
\date{Working draft, November 2026}

\begin{document}
\maketitle

\begin{abstract}
Section~2 of the paper lists how block encodings combine: products, Kronecker
products, linear combinations, stacking, concatenation and sparse matrices,
each with its subnormalization. This supplement shows, in pure matrix
terms, how each rule is realized as an explicit unitary matrix, and why the
subnormalizations are what they are. Every construction was built as a dense
matrix and checked numerically, with the script
\texttt{qp\_block\_encoding\_rules.py} in the folder \texttt{preconditioning/code}
(Table~\ref{tab:check}).
\end{abstract}

\section{Conventions}

\paragraph{Registers are Kronecker factors.} A quantum computer with $q$
qubits acts on $\C^2\otimes\cdots\otimes\C^2=\C^{2^q}$. A \emph{register} is a group of
these factors. A \emph{circuit} is a product of unitary matrices, each acting
on a few factors and as the identity on the rest. Two operations are free:
\begin{itemize}
\item permuting tensor factors, which is only a relabelling of qubits;
\item adding fresh factors that start in the first basis vector $e_0$.
\end{itemize}

\paragraph{Ordering.} We write extra registers (ancillas, index registers) as
the \emph{outer} Kronecker factors and the system as the innermost one. A
vector in $\C^2\otimes\C^n$ is therefore a stacked vector $[v_0;v_1]$, and $e_0\otimes v=[v;0]$.

\paragraph{Block encodings.} A unitary $U\in\C^{2^an\times2^an}$ is a block encoding of
$A\in\C^{n\times n}$ with subnormalization $\alpha$ if its leading block is $A/\alpha$:
\[
(e_0^T\otimes I_n)\,U\,(e_0\otimes I_n)=A/\alpha,\qquad\text{i.e.}\qquad
U=\begin{bmatrix}A/\alpha & *\\ * & *\end{bmatrix}.
\]
Applying $U$ to $[v;0]$ gives $[Av/\alpha;\ w]$ for some $w$. Measuring the
ancilla and keeping the outcome ``$0$'' (\emph{postselection}) leaves
$Av/\|Av\|$, with probability $\|Av\|^2/(\alpha\|v\|)^2$.

For rectangular $A\in\C^{m\times n}$ the rows and the columns are selected by
different projections. These are $P_{\rm out}$ and $P_{\rm in}$, both with
orthonormal columns, and we require $P_{\rm out}^*UP_{\rm in}=A/\alpha$.

\paragraph{Why $\alpha\ge\|A\|$.} Every block of a unitary matrix has norm at
most one, so $\|A/\alpha\|\le1$. The subnormalization is the price of
embedding a non-unitary matrix in a unitary one. It enters the cost of
everything downstream: success probabilities, and the degree of QSVT
polynomials.

\section{A single matrix: the unitary dilation}

\begin{lemma}[Halmos dilation]
Let $\|X\|\le1$. Then
\[
U=\begin{bmatrix} X & (I-XX^*)^{1/2}\\ (I-X^*X)^{1/2} & -X^*\end{bmatrix}
\]
is unitary.
\end{lemma}
\begin{proof}
Use $X^*f(XX^*)=f(X^*X)X^*$ for $f(t)=(1-t)^{1/2}$. This makes the off-diagonal
blocks of $U^*U$ vanish, and
\[
U^*U=\begin{bmatrix}X^*X+(I-X^*X) & 0\\ 0 & (I-XX^*)+XX^*\end{bmatrix}=I .
\]
\end{proof}

So every $A$ has a block encoding with $\alpha=\|A\|$ and one ancilla qubit. The
dilation is dense, however, and has no short circuit in general. The point
of the rules below is to build block encodings of \emph{structured} matrices
from small, cheap pieces. The price is a subnormalization that may exceed
$\|A\|$.

\section{The ideal subnormalization and what circuits achieve}\label{sec:ideal}

\paragraph{Definition.} The \emph{ideal subnormalization} of $A$ is $\alpha=\|A\|$. It is a
lower bound for every block encoding, and the Halmos dilation attains it.
The dilation, however, is dense and has no efficient circuit in general.
Efficient block encodings are assembled from structured pieces by the
rules of the following sections. At each step the subnormalization can
exceed the norm. The table below summarizes when a rule is tight, i.e.
when its $\alpha$ equals the norm of the result.

\begin{center}\small
\begin{tabular}{@{}lp{5.4cm}p{5.6cm}@{}}
\toprule
rule & $\alpha$ versus norm & tight when \\
\midrule
$AB$ & $\alpha_A\alpha_B\ge\|A\|\|B\|\ge\|AB\|$ & $A$ is maximal on the range of $B$'s top singular vector\\
$A\otimes B$ & $\alpha_A\alpha_B$ vs.\ $\|A\otimes B\|=\|A\|\|B\|$ & always (if $\alpha_A,\alpha_B$ are tight)\\
$\sum_jc_jA_j$ & $\sum_j|c_j|\alpha_j\ge\|\sum_jc_jA_j\|$ & all terms maximal on a common vector, phases aligned\\
$[A_0;\dots;A_{k-1}]$ & $(\sum_j\alpha_j^2)^{1/2}$ vs.\ $\|\sum_jA_j^*A_j\|^{1/2}$ & all $A_j$ maximal on a common input vector\\
$[A_0,\dots,A_{k-1}]$ & $(\sum_j\alpha_j^2)^{1/2}$ vs.\ $\|\sum_jA_jA_j^*\|^{1/2}$ & all $A_j$ maximal on a common output vector\\
sparse $A$ & $a_{\max}\sqrt{s_rs_c}$ (Schur-test bound) vs.\ $\|A\|$ & rarely\\
\bottomrule
\end{tabular}
\end{center}

\paragraph{The common mechanism.} A loss occurs when the pieces attain
their norms on \emph{different} vectors. In PDE discretizations the typical
conflict is between oscillatory vectors, which maximize stiffness and
gradient operators, and smooth vectors, which maximize mass and
prolongation operators.

\paragraph{Examples.}
\begin{itemize}
\item \emph{A tight LCU.} $I-S$ with the cyclic shift $S$ has $\alpha=2=\|I-S\|$
  for even $n$: both terms are maximal on the alternating vector, with
  aligned signs.
\item \emph{An LCU that loses a factor $2$.} The 2D stiffness matrix
  $K=K_1\otimes M_1+M_1\otimes K_1$ as an LCU of two Kronecker products has
  $\alpha=2\|K_1\|\|M_1\|\to8$, while $\|K\|\to4$ ($Q_1$, $h=2^{-L}$; the ratio is
  $1.993$, $2.000$, $2.000$, $2.000$ for $L=4,6,8,10$). $\|K_1\|$ is attained by the
  most oscillatory sine vector, $\|M_1\|$ by the smoothest, and no single
  vector makes both terms maximal.
\item \emph{A product that loses $2^{L-\ell}$.} $\|F_1\|\|P_\ell\|/\|F_1P_\ell\|$ (gradient factor
  times prolongation) equals $2^{L-\ell}$ up to the coarsest levels (paper,
  Table~4), for the same reason. The commuting lemma of the paper removes
  this loss by applying the gradient on the coarse level first.
\item \emph{Concatenation over levels loses $\sqrt L$.} For
  $G=[T_1G_1,\dots,T_LG_L]$ the rule gives $\alpha\approx\sqrt L\max_\ell\|G_\ell\|$. Multilevel
  theory (a strengthened Cauchy--Schwarz inequality between levels) keeps
  $\|G\|=O(1)$, because different levels mostly act on different frequency
  bands. Their output ranges are far from aligned, so the rule is not tight.
\item \emph{Sparse access loses a factor 3} in the random example of
  Table~\ref{tab:check}: $\alpha=15.3$ against $\|A\|=5.1$.
\end{itemize}

\paragraph{What this means for the paper.} The preconditioned KKT matrix is
encoded as $\B=\G^T\J\G$ with $\G=I_2\otimes G$, where
$G=[F_{\rm val}R;\beta^{1/4}F_\nabla R]$ maps multilevel coefficients to values and
gradients at quadrature points. The reflection $\J$ has $\alpha=1$, so
$\alpha_\B=\alpha_G^2$.
\begin{itemize}
\item \emph{Ideal case:} $\alpha_G=\|G\|$, giving $\alpha_\B=\|G\|^2=\|R^TA_\beta R\|\le c_1$. The paper's
  Table~1 has $\|G\|\approx1.7$--$2.7$, hence $\kappa_\alpha(\B)\approx5$--$10$.
\item \emph{Even the ideal case is above $\|\B\|$:} $\|\B\|\le\|G\|^2$ since $\|\J\|=1$. This
  is the price of writing an indefinite matrix with a reflection in the
  middle. Here both quantities are $O(1)$.
\item \emph{Actual construction:} $\|G\|\le\alpha_G\le C_d\sqrt L$, with losses from three
  places: the 1D LCUs, the stack of values and gradients, and the
  concatenation over levels. The pessimistic sparse-access constant is
  $C_3\approx45$.
\item \emph{Why measuring $\alpha_G$ matters:} it enters the cost squared. For $d=3$
  the gap between the bounds is up to $(45\sqrt L/2)^2\approx500L$. Building the
  encoding explicitly and measuring $\alpha_G$ (to-do item 1.1) decides where
  between the two bounds we are.
\end{itemize}

\section{Products: fresh ancillas}

Let $U_A=\begin{bmatrix}X_A&Y_A\\Z_A&*\end{bmatrix}$ and $U_B=\begin{bmatrix}X_B&Y_B\\Z_B&*\end{bmatrix}$
be block encodings with $X_A=A/\alpha_A$ and $X_B=B/\alpha_B$.

\paragraph{The naive product fails.} If both use the \emph{same} ancilla,
\[
U_AU_B=\begin{bmatrix}X_AX_B+Y_AZ_B & *\\ * & *\end{bmatrix}.
\]
The ``garbage'' $Z_Bv$ that $U_B$ moves into the ancilla-1 block is
brought back into the leading block by $Y_A$.

\paragraph{The construction.} Give each factor its own ancilla. Order the
space as (ancilla of $A$) $\otimes$ (ancilla of $B$) $\otimes$ system. Let $\widetilde U_A$
act as $U_A$ on the ancilla of $A$ and the system, and as the identity on the
ancilla of $B$; define $\widetilde U_B$ in the same way. Then
\[
\big(e_0\otimes e_0\otimes I\big)^T\,\widetilde U_A\widetilde U_B\,\big(e_0\otimes e_0\otimes I\big)=X_AX_B=\frac{AB}{\alpha_A\alpha_B}.
\]
The reason is that $\widetilde U_B$ puts its garbage into the block where the ancilla
of $B$ equals $1$. $\widetilde U_A$ never touches that ancilla, so the garbage
stays orthogonal to the leading block.
\begin{itemize}
\item \emph{Cost:} one call to each factor, plus the extra ancillas.
\item \emph{Subnormalization:} $\alpha_A\alpha_B$, which can be much larger than
  $\|AB\|$. This is the mechanism behind the failure of the naive product
  encoding of $\widetilde R^T\mathcal A\widetilde R$ in the paper.
\end{itemize}

\section{Kronecker products}

$U_A\otimes U_B$ acts on (ancilla$_A$, system$_A$, ancilla$_B$, system$_B$). After the
free permutation to (ancilla$_A$, ancilla$_B$, system$_A$, system$_B$), its leading
block is
\[
X_A\otimes X_B=\frac{A\otimes B}{\alpha_A\alpha_B}.
\]
No calls are combined at all: the two circuits run side by side on disjoint
qubits. This is why the Kronecker structure of $Q_1$ matrices on Cartesian
grids is so valuable. A $d$-dimensional factor costs $d$ one-dimensional
circuits.

\section{Linear combinations of unitaries (LCU)}

Let $\sum_{j=0}^{k-1}c_jA_j$ be given, with block encodings $U_j$ of $A_j$
(subnormalization $\alpha_j$), $k=2^m$ terms, and $c_j>0$. Complex phases of
$c_j$ are absorbed into $U_j$. Let $s=\sum_jc_j\alpha_j$.
\begin{itemize}
\item \emph{PREPARE:} a unitary $V\in\C^{k\times k}$ with first column
  $Ve_0=p$, where $p_j=\sqrt{c_j\alpha_j/s}$. It acts on an index register of
  $m$ qubits.
\item \emph{SELECT:} the block-diagonal unitary
  $\mathrm{SEL}=\operatorname{blkdiag}(U_0,\dots,U_{k-1})$, i.e. ``apply $U_j$ if the index
  register is $e_j$''.
\end{itemize}
Then
\[
W=(V^*\otimes I)\,\mathrm{SEL}\,(V\otimes I)
\]
has the leading block (index $0$, ancilla $0$)
\[
\sum_j\bar p_j\,p_j\,\frac{A_j}{\alpha_j}=\sum_j\frac{c_j\alpha_j}{s}\frac{A_j}{\alpha_j}=\frac1s\sum_jc_jA_j .
\]
So $\alpha=\sum_jc_j\alpha_j$.

\paragraph{Example: the difference matrix.} Let $S$ be the cyclic shift, a
permutation matrix and hence unitary, with $\alpha=1$. With $k=2$, $p=\frac1{\sqrt2}(1,1)$,
$V=H=\frac1{\sqrt2}\begin{bmatrix}1&1\\1&-1\end{bmatrix}$ and SEL $=\operatorname{blkdiag}(I,-S)$,
\[
W=(H\otimes I)\begin{bmatrix}I&0\\0&-S\end{bmatrix}(H\otimes I)=\frac12\begin{bmatrix}I-S&I+S\\I+S&I-S\end{bmatrix},
\]
a block encoding of $I-S$ with $\alpha=2$. The 1D quadrature factors $F_0$ and $F_1$
of the paper are built this way. They are two-term combinations of shifts
with diagonal weights.

\section{Stacking}

The stack $\begin{bmatrix}A_0\\ \vdots\\ A_{k-1}\end{bmatrix}\in\C^{kn\times n}$ is tall, and the index
register becomes part of the \emph{output}. Take $V$ with $Ve_0=p$, where
$p_j=\alpha_j/N$ and $N=\big(\sum_j\alpha_j^2\big)^{1/2}$, and
\[
W=\mathrm{SEL}\,(V\otimes I).
\]
On an input $e_0\otimes e_0\otimes v$ (index $0$, ancilla $0$), $V$ produces
$\sum_jp_j\,e_j\otimes e_0\otimes v$. SEL then applies $U_j$ in branch $j$. Keeping
ancilla $0$ but \emph{all} index values gives
\[
\sum_je_j\otimes\frac{p_j}{\alpha_j}A_jv=\frac1N\begin{bmatrix}A_0v\\ \vdots\\ A_{k-1}v\end{bmatrix}.
\]
This is a rectangular block encoding with $P_{\rm in}$: (index $0$, ancilla $0$), and
$P_{\rm out}$: (any index, ancilla $0$). Its subnormalization is
$N=\big(\sum_j\alpha_j^2\big)^{1/2}$. In the paper this stacks the value rows and
the $d$ gradient rows of $G_\ell$.

\section{Concatenation}

The concatenation $[A_0,\dots,A_{k-1}]\in\C^{n\times kn}$ is wide, and the index register is
part of the \emph{input}. With the same $V$,
\[
W=(V^*\otimes I)\,\mathrm{SEL}.
\]
An input $\sum_je_j\otimes e_0\otimes v_j$ goes to $\sum_je_j\otimes(e_0\otimes A_jv_j/\alpha_j+\dots)$.
Applying $V^*$ and keeping index $0$ and ancilla $0$ gives
\[
\sum_j\bar p_j\frac{A_jv_j}{\alpha_j}=\frac1N\sum_jA_jv_j .
\]
So $\alpha=N=\big(\sum_j\alpha_j^2\big)^{1/2}\le\sqrt k\max_j\alpha_j$. The factor $\sqrt k$ is
the \emph{price of summing}: the vectors $A_jv_j$ add up coherently only in
the one index branch that is kept. In the paper this is the concatenation
over the $L$ levels, $G=[T_1G_1,\dots,T_LG_L]$. It is the source of
$\alpha_G=O(\sqrt L)$, and the reason the paper asks whether the $\sqrt L$ can be
avoided.

\section{Sparse matrices}

Let $A\in\C^{n\times n}$ have at most $s_r$ nonzeros per row and $s_c$ per column, and
let $a_{\max}=\max|a_{ij}|$. Work in $\C^n\otimes\C^n\otimes\C^2$ and define, for each row
$i$ and column $j$, the unit vectors
\[
\psi_i=\frac1{\sqrt{s_r}}\sum_{j\in\mathcal R_i}e_i\otimes e_j\otimes\Big(\frac{\bar a_{ij}}{a_{\max}}e_0+\sqrt{1-\tfrac{|a_{ij}|^2}{a_{\max}^2}}\,e_1\Big),\qquad
\phi_j=\frac1{\sqrt{s_c}}\sum_{i\in\mathcal C_j}e_i\otimes e_j\otimes e_0 .
\]
Here $\mathcal R_i$ ($\mathcal C_j$) is the nonzero pattern of row $i$ (column $j$), padded
with zero entries to exactly $s_r$ ($s_c$) indices.
\begin{itemize}
\item The $\psi_i$ are orthonormal, because they differ in the first factor.
  The $\phi_j$ are orthonormal, because they differ in the second.
\item $\psi_i^*\phi_j=a_{ij}/(a_{\max}\sqrt{s_rs_c})$.
\end{itemize}
With the isometries $T_L=[\psi_1,\dots,\psi_n]$ and $T_R=[\phi_1,\dots,\phi_n]$,
\[
\frac{A}{a_{\max}\sqrt{s_rs_c}}=T_L^*T_R .
\]
Completing $T_L$ and $T_R$ to unitaries $U_L=[T_L,\,T_L^\perp]$ and
$U_R=[T_R,\,T_R^\perp]$ gives the block encoding $W=U_L^*U_R$, with
$\alpha=a_{\max}\sqrt{s_rs_c}$.

\paragraph{What ``entries computable by a circuit'' means.} $U_R$ maps $e_j$ to
$\phi_j$: it must produce the positions of the nonzeros of column $j$ (a
column-index oracle). $U_L$ maps $e_i$ to $\psi_i$: it must produce the
positions in row $i$, and rotate a flag qubit by an angle determined by
$a_{ij}$ (a value oracle). For structured matrices, such as stencils on
Cartesian grids, these are small arithmetic circuits on indices. For
unstructured meshes they need lookup tables, which is the data-access
problem noted in the paper.

\paragraph{The constant.} $\alpha=a_{\max}\sqrt{s_rs_c}$ is exactly the Schur-test bound for
$\|A\|$, so it is never smaller than $\|A\|$, and often noticeably larger. In
the random example of Table~\ref{tab:check}, $\alpha=15.3$ against $\|A\|=5.1$. This is
the origin of the pessimistic constant $C_d$ in the paper's
Proposition~3.8, and of the suggestion to use Kronecker and LCU
constructions instead of generic sparse access.

\section{Summary of costs}

\begin{table}[ht]
\centering\small
\begin{tabular}{llll}
\toprule
rule & realization & extra qubits & subnormalization\\
\midrule
single $A$ (dense) & Halmos dilation & 1 & $\|A\|$\\
$AB$ & $\widetilde U_A\widetilde U_B$, fresh ancillas & $a_A+a_B$ & $\alpha_A\alpha_B$\\
$A\otimes B$ & $U_A\otimes U_B$, permuted & $a_A+a_B$ & $\alpha_A\alpha_B$\\
$\sum_jc_jA_j$ & $(V^*\otimes I)\,\mathrm{SEL}\,(V\otimes I)$ & $\log_2k+\max a_j$ & $\sum_jc_j\alpha_j$\\
$[A_0;\dots;A_{k-1}]$ & $\mathrm{SEL}\,(V\otimes I)$ & $\log_2k+\max a_j$ & $(\sum_j\alpha_j^2)^{1/2}$\\
$[A_0,\dots,A_{k-1}]$ & $(V^*\otimes I)\,\mathrm{SEL}$ & $\log_2k+\max a_j$ & $(\sum_j\alpha_j^2)^{1/2}\le\sqrt k\max\alpha_j$\\
sparse $A$ & $U_L^*U_R$ & $\log_2n+1$ & $a_{\max}\sqrt{s_rs_c}$\\
isometry $T$ & $[T,\,T^\perp]$ & as needed & $1$\\
\bottomrule
\end{tabular}
\caption{The combination rules. SEL needs \emph{controlled} versions of the $U_j$,
which typically cost a small constant factor more gates.}
\label{tab:rules}
\end{table}

\section{\texorpdfstring{How the rules build $G$ in the paper}{How the rules build G in the paper}}

$G=[T_1G_1,\dots,T_LG_L]$ is assembled from the inside out:
\begin{enumerate}
\item 1D two-term factors (values and derivatives at Gauss points): LCUs of
  shifts with diagonal weights, $\alpha$ of the order of their norm.
\item $d$-dimensional factors: Kronecker products of 1D factors, so the
  $\alpha$'s multiply. They match the norms because the norms of Kronecker
  products also multiply.
\item $G_\ell$: a stack of the value factor and the $d$ gradient factors, weighted
  by $\omega_\ell$ and $\beta^{1/4}$, with $\alpha(G_\ell)=(\sum_c\alpha_c^2)^{1/2}$.
\item $T_\ell G_\ell$: a product with an isometry, which costs nothing in $\alpha$.
\item $G$: a concatenation over the levels, which costs $\sqrt L$.
\end{enumerate}
The only places where $\alpha$ can exceed the norm are step~1 (if the LCU is
not tight), step~3 (stacking is tight only if all components are equally
saturated) and step~5 (the $\sqrt L$). Measuring $\alpha_G$ for the explicit
construction (to-do item 1.1) quantifies exactly these three losses.

\section{Amplitude amplification: a rotation in a plane}\label{sec:aa}

A block encoding produces the wanted vector only as one block of a longer
vector. Postselection keeps that block, but succeeds only with some
probability. Amplitude amplification raises this probability to nearly one,
at a cost proportional to $1/a$ instead of $1/a^2$, where $a$ is the success
amplitude. It is the source of the factor $\kappa_\alpha^2$ in plain QSVT
inversion (Section~2 of the paper).

\paragraph{The setting.} Let $W$ be a unitary (e.g.\ the QSVT circuit for
$p(\B/\alpha)$), $b$ a unit vector, and $u=e_0\otimes b$. Let $\Pi$ be the orthogonal
projector onto the ``ancilla $=0$'' subspace, i.e.\ onto the first block. Then
\[
\psi:=Wu=\begin{bmatrix}x\\ r\end{bmatrix},\qquad x=\Pi\psi\ \text{(e.g.\ }x=\sigma_{\min}\B^+b),\qquad a:=\|x\|\le1 .
\]
Measuring the ancilla returns $x/\|x\|$ with probability $a^2$. Repeating until
success costs about $1/a^2$ applications of $W$.

\paragraph{Two reflections.} Define
\[
S_\Pi=I-2\Pi,\qquad S_\psi=I-2\psi\psi^*=W\,(I-2uu^*)\,W^*,\qquad Q:=-S_\psi S_\Pi .
\]
$S_\Pi$ flips the sign of the first block, which is a cheap circuit. $S_\psi$ needs
one $W^*$, one reflection about $u$ (which requires running the preparation
of $b$ backwards and forwards), and one $W$.

\paragraph{The invariant plane.} Let $g=\Pi\psi/a$ (the good direction) and
$h=(I-\Pi)\psi/\|(I-\Pi)\psi\|$ (the bad direction). They are orthonormal, and with
$a=\sin\theta$,
\[
\psi=\sin\theta\,g+\cos\theta\,h .
\]
Both reflections map $\operatorname{span}\{g,h\}$ into itself: $S_\Pi g=-g$, $S_\Pi h=h$, and $S_\psi$
reflects about the line orthogonal to $\psi$. In the basis $(g,h)$,
\[
S_\Pi=\begin{bmatrix}-1&0\\0&1\end{bmatrix},\qquad
-S_\psi=2\psi\psi^*-I=\begin{bmatrix}-\cos2\theta&\sin2\theta\\ \sin2\theta&\cos2\theta\end{bmatrix},\qquad
Q=\begin{bmatrix}\cos2\theta&\sin2\theta\\ -\sin2\theta&\cos2\theta\end{bmatrix},
\]
so $Q$ is the rotation by $2\theta$ towards $g$. Hence
\[
Q^k\psi=\sin\big((2k+1)\theta\big)\,g+\cos\big((2k+1)\theta\big)\,h .
\]
With $k\approx\pi/(4\theta)-\tfrac12$ the coefficient of $g$ is close to one, and a single
measurement succeeds with probability $\sin^2((2k+1)\theta)\ge1-\sin^2\theta$ for
the best integer $k$. The number of rounds is $k\approx\pi/(4a)$. Each round
costs one $W$ and one $W^*$.

\paragraph{Polynomial view.} Since $a=\sin\theta=\cos(\pi/2-\theta)$,
\[
\sin\big((2k+1)\theta\big)=(-1)^kT_{2k+1}(a),
\]
with $T_n$ the Chebyshev polynomial of the first kind. Amplitude
amplification is therefore QSVT applied to the $1\times1$ ``matrix'' $a$ (the
singular value of $\Pi W\Pi$ on $u$), with an odd polynomial of degree $2k+1$
that maps $a$ close to $\pm1$. Its degree $O(1/a)$ follows the same law as the
degree $O(1/\delta)$ of the polynomial approximating $\delta/t$ on $[\delta,1]$.

\paragraph{Application to QSVT inversion.} The polynomial acts on $\B/\alpha$, so
$p(\B/\alpha)\approx\delta(\B/\alpha)^+=\sigma_{\min}\B^+$ with $\delta=\sigma_{\min}/\alpha$. For
$x=\sigma_{\min}\B^+b$ and $b\in\operatorname{range}\B$,
\[
a=\sigma_{\min}\|\B^+b\|\ \ge\ \frac{\sigma_{\min}}{\|\B\|}\ \ge\ \frac{\sigma_{\min}}{\alpha}=\frac1{\kappa_\alpha}.
\]
So at most about $(\pi/4)\kappa_\alpha$ rounds are needed, each containing $W$ of
degree $D=O(\kappa_\alpha\log(\kappa_\alpha/\varepsilon))$: in total
$O(\kappa_\alpha^2\log(\kappa_\alpha/\varepsilon))$ calls to the block encoding of $\B$. The two
factors $\kappa_\alpha$ have different origins: one is the degree of $p$, the other
the size of $\sigma_{\min}\B^+b$. Variable-time amplitude amplification (Childs,
Kothari, Somma 2017) and the discrete adiabatic method (Costa et al.\ 2022)
avoid multiplying them and reach $O(\kappa_\alpha\log(1/\varepsilon))$.

\paragraph{Practical points.}
\begin{itemize}
\item The optimal $k$ depends on $a$, which is unknown. Exponential search
  over $k$ (Brassard, H{\o}yer, Mosca, Tapp 2002) or fixed-point amplitude
  amplification (Yoder, Low, Chuang 2014) cost only a constant factor more.
\item Too many rounds rotate past $g$; the success probability oscillates
  in $k$. This is the price of a non-adaptive polynomial, as for QSVT itself.
\item $S_\psi$ requires that $b$ can be prepared by a circuit and that this
  circuit can be inverted. This is a standing assumption of QLSAs.
\end{itemize}

\section{Removing the mass-matrix floor}\label{sec:floor}

\paragraph{Where the floor comes from.} The multilevel factor of the paper
scales each level by a scalar, $\omega_\ell^2=\theta_\ell/\delta_\ell$ with $\delta_\ell$ the (constant)
diagonal of $A_{\beta,\ell}=M_\ell+\sqrt\beta K_\ell$. On a mass-dominated level this treats
$M_\ell$ as a multiple of the identity, but $M_\ell=M_1^{\otimes d}$ has eigenvalues in
$((h/3)^d,h^d)$, so $\kappa(\delta_\ell^{-1}M_\ell)\to3^d$. When the finest grid does not resolve
$\beta$, $\kappa(CA_\beta)$ therefore approaches $3^d/(1-4^{-s})$: 4, 12 and 36 for $d=1,2,3$
and $s=1$. This is harmless for Poisson control on fine grids, but in the
time-periodic heat problem all high frequencies are mass-dominated once
$\tau\lesssim2h_L^2$, and the worst frequency sets $\kappa_\alpha$.

\paragraph{A better level scaling.} Replace the scalar $\delta_\ell^{-1}$ by
\[
D_\ell=(M_\ell+c_\ell I)^{-1},\qquad c_\ell=\sqrt\beta\,\operatorname{diag}(K_\ell)\quad(\text{a scalar}),
\]
i.e.\ the exact mass matrix plus Jacobi scaling of the stiffness part, with
$C=\sum_\ell\tilde\theta_\ell P_\ell D_\ell P_\ell^T$. On mass-dominated levels $D_\ell\approx M_\ell^{-1}$, on
stiffness-dominated levels $D_\ell\approx c_\ell^{-1}I$ as before. In the unresolved limit
$CM=\sum_\ell\tilde\theta_\ell Q_\ell$ with $Q_\ell=P_\ell M_\ell^{-1}P_\ell^TM$ the matrix of the $L^2$
projection onto $V_\ell$ ($Q_L=I$), so $\operatorname{spec}(CM)\subset[1,\sum_\ell\tilde\theta_\ell]$ and the floor
drops to $1/(1-4^{-s})=4/3$. In the proof of the paper's Theorem~3.5 the
nodal-versus-$L^2$ comparison (the factor $3d\cdot3^d$) becomes an identity on
mass-dominated levels. Two simpler variants are worse: $M_\ell^{-1}$ on
\emph{all} levels loses on resolved grids (stiffness-dominated levels want a
scalar), and switching between $M_\ell^{-1}$ and Jacobi by level type gives
$\kappa\approx10$ at intermediate $\beta$.

\paragraph{The encoding challenge.} The new level factor is
\[
G_\ell=\sqrt{\tilde\theta_\ell}\begin{bmatrix}F_{\mathrm{val}}\\ \beta^{1/4}F_\nabla\end{bmatrix}D_\ell^{1/2},\qquad
F_{\mathrm{val}}^TF_{\mathrm{val}}=M_\ell,\quad F_\nabla^TF_\nabla=K_\ell .
\]
$M_\ell$, $K_\ell$ and $D_\ell$ are all diagonalized by the $d$-dimensional sine transform,
so everything can be read off mode by mode, with $\mu$ and $k$ the eigenvalues
of $M_\ell$ and $K_\ell$.
\begin{itemize}
\item \emph{Value rows: (W2) fails for the plain product.} $F_{\mathrm{val}}$ is largest on
  smooth modes ($\mu\approx h^d$), $D_\ell^{1/2}$ on oscillatory ones ($\mu\approx(h/3)^d$). On a
  mass-dominated level $\|F_{\mathrm{val}}\|\,\|D_\ell^{1/2}\|\approx3^{d/2}$, whereas
  $\|F_{\mathrm{val}}D_\ell^{1/2}\|\le1$. Encoding the product of the two factors
  therefore inflates $\alpha$ by $3^{d/2}$ and $\kappa_\alpha$ by $3^d$: the floor moves from
  $\kappa(CA_\beta)$ into the subnormalization.
\item \emph{Gradient rows: nearly aligned.} $F_\nabla$ and $D_\ell^{1/2}$ are both largest on
  oscillatory modes; the loss is a bounded factor.
\end{itemize}

\paragraph{The remedy: factor through an isometry.} Since $F_{\mathrm{val}}^TF_{\mathrm{val}}=M_\ell$
exactly, $U_{\mathrm{val}}:=F_{\mathrm{val}}M_\ell^{-1/2}$ is an isometry, and
\[
F_{\mathrm{val}}D_\ell^{1/2}=U_{\mathrm{val}}\,g(M_\ell),\qquad g(\mu)=\sqrt{\mu/(\mu+c_\ell)}\in[0,1].
\]
A product with an isometry is always tight. Both factors can be encoded with
$\alpha$ equal to their norm, using the construction of Yang, Yu and Zhang
(arXiv:2610.06987, Secs.~3.4--3.5):
\begin{enumerate}
\item \emph{Functions of $M_\ell$} (here $g(M_\ell)$, there $M_\ell^{-1/2}$) are applied in
  the sine basis: a sine transform (via the QFT), a rotation of a flag
  qubit by an angle computed from the mode index, so that the amplitude on
  the flag state $0$ is $g(\mu_j)/\max g$, and the inverse transform. This is a
  block encoding with $\alpha=\max g$, i.e.\ tight; no polynomial
  approximation is needed. For $D_\ell$ the angle depends on all $d$ mode
  indices, since $M_\ell+c_\ell I$ is not a Kronecker product; the arithmetic is
  still cheap.
\item \emph{The isometry $U_{\mathrm{val}}$} is the Kronecker product of $d$ one-dimensional
  isometries $E_1M_1^{-1/2}$ ($E_1$: values at the two Gauss points). The product
  encoding of $E_1$ and $M_1^{-1/2}$ has $\alpha=\sqrt{4h/3}\,\sqrt{3/h}=2$, i.e.\ its block is
  $F=\tfrac12E_1M_1^{-1/2}$ with $F^*F=\tfrac14I$. One step of \emph{oblivious} amplitude
  amplification (Section~\ref{sec:aa}, with the reflection about the input
  replaced by the reflection about the ancilla state $0$) has block
  $3F-4FF^*F=2F$, i.e.\ it encodes the isometry with $\alpha=1$, using three calls
  to $U$ or $U^*$. This works because $F^*F$ is a multiple of the identity:
  all singular values are rotated by the same angle. The $d$-dimensional
  $U_{\mathrm{val}}$ then has $\alpha=1^d=1$.
\item \emph{Gradient rows.} The same idea works per direction $q$: with
  $E_1=U_1M_1^{1/2}$,
  \[
  F_\nabla^{(q)}D_\ell^{1/2}=\big[U_1\otimes\dots\otimes C_1\otimes\dots\otimes U_1\big]\cdot
  \Big[\big(M_1^{1/2}\otimes\dots\otimes I\otimes\dots\otimes M_1^{1/2}\big)D_\ell^{1/2}\Big],
  \]
  where $C_1$ (derivatives, $\alpha=\|C_1\|=2/\sqrt h$) sits in position $q$. The first
  factor is a Kronecker product of isometries and $C_1$, hence tight; the
  second is diagonal in the sine basis, hence tight. The two are nearly
  aligned (both favour modes that oscillate in direction $q$), so the
  product loses only a bounded factor. The directions are stacked.
\end{enumerate}

\paragraph{Measured effect.} Table~\ref{tab:floor} compares, with the same
one-dimensional primitives for all three: the paper's Jacobi scaling, the
Yang--Yu--Zhang factor ($M_\ell^{-1/2}$ on every level, weights $(1+t_\ell)^{-1}$, i.e.\
$s=0$), and the hybrid scaling $D_\ell$ with $s=1$. The subnormalization is
$\alpha_G^2=\sum_\ell\alpha_\ell^2$ (concatenation) and $\kappa_\alpha=\alpha_G^2/\sigma_{\min}(\B)$.

\begin{table}[ht]
\centering\small
\begin{tabular}{rrr|rrr|rrrr}
\toprule
 & & & \multicolumn{3}{c|}{$\kappa(CA_\beta)$} & \multicolumn{4}{c}{$\kappa_\alpha=\alpha_G^2/\sigma_{\min}(\B)$}\\
$d$ & $L$ & $\beta$ & Jacobi & YYZ & hybrid & Jacobi, raw & Jacobi, tight & YYZ & hybrid, tight\\
\midrule
1 & 8 & 1 & 5.0 & 7.3 & 5.0 & 16.0 & 16.0 & 35.7 & 16.0\\
1 & 8 & $10^{-4}$ & 5.4 & 7.2 & 6.2 & 17.2 & 16.7 & 35.4 & 18.3\\
1 & 8 & $10^{-8}$ & 3.6 & 4.4 & 4.5 & 7.8 & 7.4 & 17.1 & 8.3\\
1 & 8 & $10^{-12}$ & 2.3 & 7.0 & 1.3 & 5.0 & 3.9 & 9.3 & 2.1\\
\midrule
2 & 5 & 1 & 4.1 & 9.8 & 4.1 & 26.6 & 18.0 & 37.8 & 18.1\\
2 & 5 & $10^{-4}$ & 2.9 & 8.9 & 3.6 & 19.0 & 13.6 & 29.8 & 15.6\\
2 & 5 & $10^{-8}$ & 3.8 & 4.1 & 1.3 & 12.9 & 8.0 & 8.3 & 3.6\\
2 & 5 & $10^{-12}$ & 11.5 & 5.0 & 1.3 & 21.2 & 11.8 & 5.0 & 1.4\\
\midrule
3 & 3 & 1 & 4.2 & 9.7 & 4.1 & 43.2 & 17.8 & 30.2 & 17.8\\
3 & 3 & $10^{-4}$ & 2.8 & 6.9 & 2.4 & 22.1 & 10.5 & 16.0 & 8.6\\
3 & 3 & $10^{-8}$ & 20.5 & 3.0 & 1.5 & 69.8 & 25.8 & 3.3 & 2.0\\
3 & 3 & $10^{-12}$ & 24.4 & 3.0 & 1.3 & 67.4 & 24.5 & 3.0 & 1.3\\
\bottomrule
\end{tabular}
\caption{Preconditioner quality and effective condition number for three
level scalings. ``Raw'': the one-dimensional value factor $E_1$ encoded
directly ($\alpha=\sqrt{4h/3}$ against $\|E_1\|=\sqrt h$). ``Tight'': $E_1=U_1M_1^{1/2}$ with the
isometry $U_1$ and a diagonal factor in the sine basis. $\alpha_G$ is computed from
the combination rules with these primitive subnormalizations
(\texttt{qp\_mass\_floor.py}).}
\label{tab:floor}
\end{table}

Three observations.
\begin{itemize}
\item \emph{The remedy works.} In the unresolved 3D cases it reduces
  $\kappa_\alpha$ from about $25$ (Jacobi, tight) to $1.3$. On resolved grids it is within
  about $15\%$ of the Jacobi route, sometimes slightly better, sometimes
  slightly worse. The floor does not reappear in $\alpha$.
\item \emph{Tight primitives help the paper's construction too.} Even with
  Jacobi scaling, encoding the value factor through the isometry $U_1$ lowers
  $\kappa_\alpha$ by up to a factor $2.4$ in 3D (e.g.\ $43\to18$ at $\beta=1$), because the
  raw value factor loses $(2/\sqrt3)^d$ in $\alpha^2$ per level and the raw gradient
  factor more. This is relevant for Proposition~3.8 and to-do item~1.1.
\item \emph{Yang--Yu--Zhang} encode tightly but do not damp mass-dominated
  levels and use $M_\ell^{-1}$ also on stiffness-dominated levels; their
  $\kappa(CA_\beta)$ is up to $9.8$ on resolved grids, and their $\kappa_\alpha$ is about twice
  ours in 1D and 2D.
\end{itemize}

\paragraph{What remains.}
\begin{itemize}
\item The gate cost per call to $G$: per level a sine transform and its
  inverse (QFT on $\ell+1$ qubits per dimension), reversible arithmetic for the
  rotation angle, and the three-call amplification of the value isometry.
  All of this is polylogarithmic, but it multiplies the cost of every call
  to $G$ by a constant that has to be weighed against the gain in $\kappa_\alpha$.
\item Adapting the robustness proof to $D_\ell$ (straightforward: the
  nodal lemma becomes an identity on mass-dominated levels).
\item Prior art: the encoding of $M_\ell^{-1/2}$ and of the normalized value map
  is due to Yang--Yu--Zhang. New would be the hybrid scaling $D_\ell$ with damped
  weights, which is what makes the floor disappear without losing on
  resolved grids.
\end{itemize}

\section{Numerical verification}

\begin{table}[ht]
\centering\small
\begin{tabular}{lcc}
\toprule
construction & $\|W^*W-I\|_{\max}$ & error of the selected block\\
\midrule
single matrix (Halmos) & $1.2\cdot10^{-15}$ & $0$\\
product, fresh ancillas & $1.9\cdot10^{-15}$ & $8.4\cdot10^{-17}$\\
product, \emph{same} ancilla & -- & $7.9\cdot10^{-1}$ (fails, as predicted)\\
Kronecker product & $1.3\cdot10^{-15}$ & $5.6\cdot10^{-17}$\\
LCU, $k=4$ complex coefficients & $3.3\cdot10^{-15}$ & $3.9\cdot10^{-16}$\\
stacking, $k=4$ & $1.3\cdot10^{-15}$ & $3.9\cdot10^{-17}$\\
concatenation, $k=4$ & $1.6\cdot10^{-15}$ & $3.9\cdot10^{-17}$\\
sparse, $n=8$, $s_r=s_c=4$ & $1.7\cdot10^{-15}$ & $0$\\
amplitude amplification, $a=0.23$, 3 rounds & -- & $2.5\cdot10^{-15}$ (rotation formula)\\
\bottomrule
\end{tabular}
\caption{Dense checks of all constructions
(\texttt{qp\_block\_encoding\_rules.py}, random complex matrices, $n=4$ unless
stated).}
\label{tab:check}
\end{table}

\end{document}
